\subsection{主理想整环的定义}

回顾（I，第 104 页），交换环 A 的一个理想称为主理想，如果它具有形式 \((a) = A a\) ，对某个 \(a \in A\) .

定义1. --- 主理想整环是一个整环（I，第116页），其中每个理想都是主理想。

例子. --- 有理整数环 Z 是一个主理想整环（I，第 111 页）。若 K 是一个交换域，则 K 上一个未定元的多项式环 K{[}X{]} 是主理想整环（IV，第 11 页，命题 11）；形式幂级数环 K{[}{[}X{]}{]} 同样如此，因为该环的每个理想都具有形式 (X \(^{n}\) ) （IV，第 38 页，命题 12）。* p 进域中的整数环是一个主理想整环。*

如果 \(\mathbf{Q}(i)\) 表示由有理数域 Q 添加不可约多项式 \(X^{2}+1\) 的一个根 i 所得到的域，则形如 \(a+bi\) 的元素（ \(\mathbf{Q}(i)\) 中的，其中 a 和 b 是有理整数）构成 \(\mathbf{Q}(i)\) 的一个子环 A，称为「高斯整数」，它是一个主理想整环（VII，第 50 页，习题 7）。相比之下，在域 \(\mathbf{Q}(\rho)\) （其中 \(\rho\) 是 \(X^{2}+5\) 的一个根）中，由元素 \(a+b\rho\) （其中 a 和 b 为有理整数）构成的子环 B 不是主理想整环（VII，第 51 页，习题 12）。

域 K 上的两个未定元的多项式环 K{[}X, Y{]} 不是主理想整环。事实上，非零常数是仅有的同时整除 X 和 Y 的元素，而它们中没有一个能生成由 X 和 Y 生成的理想。

